\chapter{Uterine disorders}
\index{Uterine disorders}

\section*{Definition and Overview}
The uterus, or womb, is a hollow, pear-shaped muscular organ in the female pelvis where a pregnancy develops. Disorders of the uterus are very common and range from benign growths that cause heavy or painful periods, to infections, structural abnormalities, and cancers. The uterus plays a central role in menstruation, fertility, and childbirth, so disorders of the organ commonly present with changes in bleeding, pain, or the ability to conceive. The most frequently encountered uterine disorders include uterine fibroids (leiomyomas), endometriosis, adenomyosis, endometrial polyps, endometrial hyperplasia, and endometrial cancer. Uterine prolapse, in which the uterus drops down into the vagina, is another important condition, particularly in older women. Rarer structural disorders, such as congenital uterine anomalies (for example, a septate or bicornuate uterus) and intrauterine adhesions (Asherman syndrome), also fall under this heading and are important causes of pregnancy problems. Because these conditions share many symptoms, such as abnormal bleeding, pelvic pain, and pressure, a clear diagnosis requires a careful history, examination, and often ultrasound. Understanding the type of disorder matters, because treatments differ widely, from reassurance and simple analgesia to hormonal treatment, minor procedures, or surgery.

\section*{Epidemiology}
Uterine disorders are among the most common reasons for gynaecological consultation. Uterine fibroids are the most common benign tumour of the female reproductive tract, affecting up to 70--80\% of women of African ancestry and 30--40\% of white women by the time of menopause. Endometriosis affects roughly 1 in 10 women of reproductive age worldwide, although many cases go undiagnosed. Adenomyosis, in which endometrial tissue grows into the muscle wall of the uterus, is found in up to 20\% of women having hysterectomies. Endometrial polyps are found in about 10--25\% of women investigated for abnormal bleeding. Endometrial cancer is the most common gynaecological cancer in developed countries, with the majority of cases occurring after menopause, usually between 60 and 70 years of age. Uterine prolapse becomes more common with age and after vaginal childbirth, affecting a significant minority of older women. These figures probably underestimate the true burden of disease, because many women do not seek care for symptoms they consider normal, and conditions such as endometriosis can only be confirmed surgically. Awareness of these conditions is growing, and more women are now presenting earlier with heavy periods, pelvic pain, and fertility concerns.

\section*{Aetiology and Pathophysiology}
The uterus has three layers: the inner lining (endometrium), a muscular wall (myometrium), and an outer serosal coat. Each layer gives rise to characteristic disorders.

\begin{itemize}
    \item \textbf{Fibroids (leiomyomas):} benign smooth-muscle tumours of the myometrium whose growth is driven by oestrogen and progesterone; they tend to shrink after menopause
    \item \textbf{Endometriosis:} endometrial-like tissue growing outside the uterus, most commonly on the ovaries, fallopian tubes, and pelvic peritoneum, causing inflammation, scarring, and pain; the exact mechanism is not fully understood, with retrograde menstruation the leading theory
    \item \textbf{Adenomyosis:} endometrial tissue within the myometrium, causing a diffusely enlarged, tender, heavy-bleeding uterus, often associated with increasing age and prior childbirth
    \item \textbf{Endometrial polyps:} benign overgrowths of endometrial glandular tissue attached to the lining by a stalk; they are hormone-responsive and more common in perimenopause
    \item \textbf{Endometrial hyperplasia:} thickening of the endometrial lining driven by unopposed oestrogen, which can, in some forms, progress to endometrial cancer
    \item \textbf{Endometrial cancer:} usually adenocarcinoma arising from the endometrium; risk factors include increasing age, obesity, diabetes, unopposed oestrogen, tamoxifen use, and Lynch syndrome
    \item \textbf{Uterine prolapse:} weakening of the supporting ligaments and pelvic floor, often after childbirth, heavy lifting, or with age, allowing the uterus to descend into the vaginal canal
    \item \textbf{Congenital uterine anomalies:} the uterus forms from two tubes that usually fuse; failure of fusion produces a bicornuate (heart-shaped) uterus, and failure of the central wall to resorb produces a septate uterus, both of which can affect pregnancy
    \item \textbf{Intrauterine adhesions (Asherman syndrome):} scarring within the uterine cavity, usually after surgery such as curettage, that can cause scanty periods or amenorrhoea and difficulty conceiving
\end{itemize}

\section*{Clinical Features}
The symptoms of uterine disorders overlap considerably, which is why diagnosis can be challenging.

\begin{itemize}
    \item Heavy or prolonged menstrual bleeding (menorrhagia), flooding, and large clots, seen with fibroids, adenomyosis, polyps, and hyperplasia
    \item Painful periods (dysmenorrhoea) that worsen over time, typical of endometriosis and adenomyosis
    \item Pelvic pain that is not confined to menstruation, including deep pain during intercourse
    \item Bleeding between periods, after intercourse, or after menopause
    \item A sensation of pelvic pressure, a lump at the vaginal opening, or a dragging feeling, suggesting prolapse
    \item Difficulty emptying the bladder or bowel, or recurrent urine infections, in advanced prolapse
    \item Subfertility or difficulty conceiving, particularly with fibroids, endometriosis, or uterine structural abnormalities
    \item Recurrent miscarriage, which is a feature of congenital uterine anomalies, intrauterine adhesions, and large submucous fibroids
    \item Abnormal discharge, fever, or general unwellness, which may indicate infection
\end{itemize}

\section*{Differential Diagnosis}
\begin{itemize}
    \item Pregnancy-related problems, including miscarriage and ectopic pregnancy
    \item Uterine fibroids, polyps, adenomyosis, and endometriosis
    \item Endometrial hyperplasia and endometrial cancer
    \item Infection of the uterus, cervix, or pelvis (endometritis, pelvic inflammatory disease)
    \item Ovarian cysts and ovarian cancer, which can cause similar pelvic pain and pressure
    \item Cervical conditions, including cervical polyps and cervical cancer
    \item Coagulation or bleeding disorders causing heavy periods
    \item Uterine prolapse and pelvic organ prolapse involving the bladder or bowel
\end{itemize}

\section*{When to Seek Medical Care}
A doctor should be consulted for heavy, painful, or irregular bleeding, bleeding after menopause, pelvic pain, bleeding between periods or after sex, or any mass or pressure symptoms that persist or worsen.

\textbf{Contact your local emergency services for:}
\begin{itemize}
    \item Heavy vaginal bleeding with dizziness, fainting, or signs of shock
    \item Severe acute pelvic pain, especially one-sided, which may indicate a ruptured cyst, ectopic pregnancy, or torsion
    \item Fever with pelvic pain and discharge, which may indicate serious pelvic infection
    \item Sudden onset of severe bleeding that soaks a pad within an hour
\end{itemize}

\section*{Diagnosis and Investigations}
A structured approach identifies the exact disorder, since management depends on it.

\begin{itemize}
    \item \textbf{History:} menstrual pattern, bleeding volume, pain, sexual history, pregnancies, medications, and family history of gynaecological or bowel cancer
    \item \textbf{Pregnancy test:} always considered first in women of reproductive age
    \item \textbf{Pelvic examination:} to assess uterine size, position, tenderness, and descent, and to detect polyps or masses at the cervix
    \item \textbf{Pelvic ultrasound:} the key initial imaging test; it identifies fibroids, adenomyosis, polyps, endometrial thickness, and ovarian findings, and can be performed by the abdomen or transvaginally for better detail
    \item \textbf{Hysteroscopy:} a camera passed into the uterine cavity to view the lining directly and to remove or biopsy polyps and small fibroids
    \item \textbf{Endometrial biopsy:} sampling of the lining, essential in women over 40, in those with risk factors, or when the lining is thickened on ultrasound
    \item \textbf{Blood tests:} full blood count to assess for anaemia from blood loss, iron studies if anaemia is found, and tumour marker CA125 only where clinically indicated
    \item \textbf{Laparoscopy:} keyhole surgery that both diagnoses and treats endometriosis
    \item \textbf{MRI:} useful in complex cases, including deep endometriosis or planning fibroid surgery
\end{itemize}

\section*{Management}
Treatment depends on the disorder, the severity of symptoms, the woman's age, and whether she wishes to have children.

\begin{itemize}
    \item \textbf{Heavy bleeding without an obvious structural cause:} tranexamic acid and anti-inflammatory medicines (NSAIDs) during menstruation, the levonorgestrel-releasing intrauterine system (Mirena), or combined oral contraception
    \item \textbf{Fibroids:} hormonal treatments to shrink them, or surgical options including myomectomy (removal of fibroids, preserving the uterus), uterine artery embolisation, and hysterectomy
    \item \textbf{Endometriosis and adenomyosis:} analgesia, hormonal treatments including the combined pill, progestogens, and GnRH analogues, laparoscopic excision of endometriosis, and in severe cases hysterectomy for adenomyosis
    \item \textbf{Endometrial polyps:} removal by hysteroscopic polypectomy
    \item \textbf{Endometrial hyperplasia:} progestogen therapy or the Mirena system to reverse the overgrowth; hysterectomy may be advised for the atypical forms
    \item \textbf{Endometrial cancer:} usually treated with hysterectomy and removal of the fallopian tubes and ovaries, sometimes with radiotherapy or chemotherapy, and with a good outlook when detected early
    \item \textbf{Uterine prolapse:} pelvic floor exercises, a vaginal pessary, or surgical repair
    \item \textbf{Congenital uterine anomalies:} surgical correction of a septum (hysteroscopic septoplasty) can improve the outcome of future pregnancies
    \item \textbf{Intrauterine adhesions:} division of the adhesions by hysteroscopy, usually followed by a course of oestrogen to regrow the lining
    \item \textbf{Fertility concerns:} surgical treatment of fibroids, endometriosis, and structural abnormalities can improve the chances of conception, and assisted reproduction is an option where needed
\end{itemize}

\section*{Complications}
\begin{itemize}
    \item \textbf{Complications:} anaemia and fatigue from chronic heavy bleeding
    \item Chronic pelvic pain that affects work, sleep, and relationships
    \item Painful intercourse, which is common with endometriosis and adenomyosis and is often unmentioned by women unless asked directly
    \item Subfertility and infertility
    \item Progressive endometriosis causing scarring and adhesions
    \item Progression of atypical endometrial hyperplasia to endometrial cancer
    \item Kidney obstruction or urinary complications in severe prolapse, which can occur when the ureters are compressed
    \item Psychological distress and reduced quality of life
\end{itemize}

\section*{Prognosis and Outcomes}
The outlook is generally excellent for benign uterine disorders, and many are manageable without surgery. Fibroids typically stop growing after menopause, and their symptoms usually resolve. Endometriosis and adenomyosis are chronic conditions that can be well controlled but may recur after treatment stops. Endometrial hyperplasia often resolves with hormonal treatment, especially if caught early. Endometrial cancer has a good prognosis when diagnosed early, and most cases are detected at an early stage because they cause bleeding. Prolapse can be improved greatly with conservative measures or surgery. Congenital uterine anomalies are often treated effectively with surgery, and many women with a bicornuate or septate uterus have successful pregnancies. Asherman syndrome responds well to surgical treatment in many cases, although severe scarring can be harder to treat. Regular follow-up is important for any condition that involves the uterine lining, because changes can develop over time.

\section*{Prevention and Self-Care}
\begin{itemize}
    \item Track your periods and report heavy, painful, or irregular bleeding early
    \item Maintain a healthy weight, since obesity increases the risk of endometrial hyperplasia, endometrial cancer, and fibroids
    \item Keep physically active and eat a balanced diet
    \item Manage pain with heat, rest, and simple anti-inflammatory medicines as advised
    \item Have routine gynaecological reviews and attend recommended screening for cervical cancer
    \item Do pelvic floor exercises to reduce the risk of prolapse, especially after childbirth
    \item Seek early review for bleeding after menopause or bleeding between periods
    \item Report any family history of uterine, bowel, or ovarian cancer to your doctor
    \item Take iron supplements if a blood test shows iron deficiency from heavy bleeding, and eat an iron-rich diet
    \item Do not accept heavy or painful periods as simply something to be endured; effective treatments are available
\end{itemize}

\section*{Sources and Further Reading}
The following organisations provide reliable information on disorders of the uterus.

\begin{itemize}
    \item Jean Hailes for Women's Health. \textit{Uterine health and heavy periods.} https://www.jeanhailes.org.au/
    \item NHS. \textit{Fibroids, endometriosis, and abnormal uterine bleeding.} https://www.nhs.uk/
    \item Mayo Clinic. \textit{Uterine fibroids, endometriosis, and endometrial cancer.} https://www.mayoclinic.org/
    \item MSD Manual. \textit{Uterine and pelvic disorders.} https://www.msdmanuals.com/
    \item MedlinePlus. \textit{Uterine diseases and conditions.} https://medlineplus.gov/
    \item Cancer Council Australia. \textit{Uterine cancer information.} https://www.cancer.org.au/
\end{itemize}
